\subsection{Measured time to first answer audio}
\label{app:ttfa}

\paragraph{Clock and endpoints.}
The serving experiment uses the same 1,190 query IDs across five pools
and five arms, with one attempt per query and arm. Input chunks become
available at their scheduled arrival times under real-time replay;
consumption cannot precede arrival. A worker's monotonic clock measures
\[
  \mathrm{TTFA}_{\rm server}
  = t_{\rm first\ answer\ PCM} - t_{\rm input\ end},
\]
where input end is the scheduled arrival of the last real audio sample.
On the local path, the endpoint is the return of the first non-listen
chunk containing nonempty PCM. This interface is chunk-blocking. On the
expert path, the endpoint is the first nonempty PCM yielded by the
teacher-forced talker, before whole-waveform synthesis completes.
The onset chunk's candidate audio and subsequent stall audio are logged
separately and excluded from answer TTFA. Negative times are retained:
the model may answer or finish before all input arrives. The measurement
does not include network delivery or client playback and does not certify
answer correctness or full-input coverage.

\paragraph{Serving configuration and calibration scope.}
Audio generation is enabled on both paths, with top-$k=20$, three forced
listen chunks, and an L22 read aggregating eight tail states and the
input mean. At escalation, ASR receives only the consumed real audio
prefix; \texttt{gpt-5.5} uses low reasoning effort, with web-enabled
requests and a non-web fallback. Prepared expert text is relayed by the
teacher-forced talker. After the local turn, expert completion is polled
at 50\,ms intervals with a 150\,s wait limit. Each worker warms up with
six chunks and one synthesis before processing sessions serially.

The frozen timing gate was fitted on 5,228 rows, with English thresholds
0.9057/0.7973/0.5534. It differs from the 5,221-row overlap-remediated fit
described for content evaluation in Appendix~\ref{app:setup-details}.
The seven external-question overlaps remain in this timing gate's
calibration. These timings therefore characterize this serving
configuration, not the generalization or joint accuracy--latency
performance of the overlap-remediated gate. No correctness judgments
were collected for the timing sessions. Hardware logs contain 5,719
sessions on H100 80GB HBM3, 56 on H100 NVL, and 175 on H200; all 1,200
Internal sessions use H100 80GB HBM3. External timings include this
hardware variation and are not controlled hardware comparisons.

\paragraph{Coverage, failures, and tails.}
All 5,950 attempts are retained: 5,904 completed sessions, 44 without a
speak commit, one empty response, and one expert timeout. Smoke runs are
excluded. Table~\ref{tab:latency} conditions timing statistics on completed
sessions with a recorded answer endpoint and reports failed attempts
separately. Calls is the fraction of all attempted sessions that fired;
this denominator includes sessions without onset. No timestamp is
zero-filled and no negative value or outlier is clipped. Internal local
TTFA has mean $-0.61$\,s but P50 0.70\,s, with 17 early responses.
Speech Web Questions has long delayed-onset outliers: local TTFA reaches
305.52\,s despite P95 3.60\,s. Quantiles describe variation across queries,
not uncertainty over repeated runs.

On the 235 Internal IDs valid in all five arms, paired-subset P50 values
are 0.69/0.88/1.61/2.88/6.64\,s for local/conservative/balanced/aggressive/
always. The corresponding common-ID counts for Speech TriviaQA, Speech
Web Questions, Speech Llama Questions, and SD-QA are 248/237/242/199.
The release records per-session timestamps, status, input identity,
expert-prefix identity, audio metadata, and frozen configuration hashes.

\begin{table}[!ht]
\centering\small
\caption{Measured server PCM-ready TTFA (seconds). Each arm attempts every query in the pool. Statistics use completed sessions; Fail counts all other outcomes, and Early counts negative TTFA among completed sessions. Calls uses all attempted sessions as its denominator. One run per query and arm; quantiles describe query variation.}
\label{tab:latency}
\setlength{\tabcolsep}{5pt}
\begin{tabular}{lrrrrrrrr}
\toprule
arm & Valid & Fail & Early & Calls (\%) & Mean & P50 & P95 & P99 \\
\midrule
\multicolumn{9}{l}{\textit{Internal ($n=240$ per arm)}} \\
local & 239 & 1 & 17 & 0.0 & -0.61 & 0.70 & 2.62 & 3.64 \\
conservative & 238 & 2 & 12 & 9.6 & 1.53 & 0.89 & 10.85 & 34.08 \\
balanced & 236 & 4 & 11 & 25.4 & 2.36 & 1.61 & 14.41 & 26.91 \\
aggressive & 236 & 4 & 13 & 50.8 & 4.59 & 2.87 & 18.14 & 54.78 \\
always & 237 & 3 & 7 & 99.2 & 8.01 & 6.67 & 25.86 & 60.83 \\
\midrule
\multicolumn{9}{l}{\textit{Speech TriviaQA ($n=250$ per arm)}} \\
local & 250 & 0 & 1 & 0.0 & 1.14 & 0.90 & 1.92 & 2.71 \\
conservative & 248 & 2 & 0 & 3.6 & 1.55 & 0.98 & 2.63 & 14.88 \\
balanced & 250 & 0 & 0 & 18.4 & 2.45 & 1.46 & 8.46 & 23.74 \\
aggressive & 250 & 0 & 0 & 56.4 & 3.62 & 3.55 & 7.92 & 16.07 \\
always & 250 & 0 & 0 & 100.0 & 5.71 & 4.89 & 10.30 & 21.76 \\
\midrule
\multicolumn{9}{l}{\textit{Speech Web Questions ($n=250$ per arm)}} \\
local & 246 & 4 & 0 & 0.0 & 3.69 & 1.71 & 3.60 & 4.25 \\
conservative & 246 & 4 & 0 & 5.6 & 5.14 & 1.81 & 7.71 & 80.26 \\
balanced & 244 & 6 & 0 & 21.2 & 3.62 & 2.18 & 9.88 & 20.61 \\
aggressive & 247 & 3 & 0 & 59.6 & 7.40 & 5.07 & 13.25 & 24.50 \\
always & 248 & 2 & 0 & 99.2 & 7.49 & 6.13 & 16.71 & 23.13 \\
\midrule
\multicolumn{9}{l}{\textit{Speech Llama Questions ($n=250$ per arm)}} \\
local & 249 & 1 & 0 & 0.0 & 1.61 & 1.52 & 2.55 & 3.66 \\
conservative & 249 & 1 & 0 & 0.0 & 1.63 & 1.54 & 2.55 & 3.71 \\
balanced & 247 & 3 & 0 & 4.0 & 2.10 & 1.53 & 3.82 & 22.78 \\
aggressive & 246 & 4 & 0 & 19.2 & 2.61 & 1.68 & 8.49 & 13.21 \\
always & 249 & 1 & 0 & 99.6 & 5.13 & 4.32 & 9.45 & 15.69 \\
\midrule
\multicolumn{9}{l}{\textit{SD-QA ($n=200$ per arm)}} \\
local & 200 & 0 & 1 & 0.0 & 0.93 & 0.62 & 1.67 & 2.24 \\
conservative & 199 & 1 & 2 & 6.5 & 1.50 & 0.75 & 7.18 & 11.94 \\
balanced & 200 & 0 & 0 & 24.5 & 2.60 & 0.87 & 10.81 & 15.52 \\
aggressive & 200 & 0 & 0 & 65.0 & 4.66 & 4.20 & 10.17 & 14.52 \\
always & 200 & 0 & 0 & 100.0 & 5.85 & 5.10 & 10.64 & 14.00 \\
\bottomrule
\end{tabular}
\end{table}

